\chapter{Homework 5 (due 10/2)}


\subsection{Revisiting the quantum Zeno effect}

We will now revisit the quantum Zeno effect in the lens of generalized measurements. Consider an initial state of a single qubit 
\begin{align}\label{quantum zeno:initial state}
\hat{\rho}_0=\dket{+}\dbra{+},~~~\dket{+}=\frac{\dket{0}+\dket{1}}{\sqrt2}
\end{align}
evolving under the following Hamiltonian
\begin{align}
\hat H=-\mu_BB\hat Z,~~~T=\frac{\pi m_e}{eB}
\end{align}
where $\mu_B=\hbar e/2m_e$ is the Bohr magneton. After a small period $\Delta t=T/n$ of unitary time evolution, we carry out a projective measurement on Pauli matrix $\hat X$. Recall that while $\dket{0},\dket{1}$ are eigenstates of Pauli $\hat Z$ operator, $\dket{\pm}$ are eigenstates of $\hat X$.  

(1) As we have learnt in the lecture, the 2-step procedure, of unitary time evolution for $\Delta t=T/n$ followed by a projective measurement on $\hat X$, can be described by a generalized measurement. Compute the measurement operators $\{\hat M_m|m=\pm\}$ as $2\times2$ matrices, and show that they satisfy the completeness relation
\begin{align}
\sum_{m=\pm}\hat M^\dagger_m\hat M_m=\hat1. 
\end{align}

(2) In the generalized measurement discussed above, if we do not know the measurement outcomes, the density matrix after the measurement is given by the action of a quantum channel $\mathcal{E}$, which is characterized by measurement operators $\{\hat M_m\}$:
\begin{align}\label{quantum zeno:quantum channel}
\mathcal{E}[\hat\rho]\equiv\sum_m\hat M_m\hat\rho\hat M_m^\dagger
\end{align}
We have previously learnt that any single-qubit density matrix can be expressed in terms of the Bloch ball:
\begin{align}
\hat\rho=\frac{\hat 1+\vec r\cdot\vec\sigma}2,~~~|\vec r|\leq1.
\end{align}
How does the vector $\vec r$ transform under the quantum channel (\ref{quantum zeno:quantum channel})?

(3) In the discussion of quantum Zeno effect, such a generalized measurement (without knowing measurement outcomes) was performed $n$ times. Starting from the initial density matrix $\hat\rho_0$ in (\ref{quantum zeno:initial state}), what is the final density matrix $\mathcal{E}^n[\rho_0]$ after $n$ such operation? 

Hint: use the results of part (2). 


(4) Using the above results, compute the probability at time $T$ (after $n$ such 2-step procedures) to produce the measurement outcome of $\sigma_x=+1$. 

Comment: recall that we have obtained this probability in HW2. You should compare your results here to HW2.  

\subsection{Measurements}

(1) \textbf{Bonus: distinguishing states by POVM measurements} \nielsen~Exercise 2.64. 

(2) Let $\dket{\psi=a\dket{0}+b\dket{1}}$ be any single-qubit quantum state. Suppose we receive one of the four states $\dket\psi,~\hat X\dket\psi,~\hat Y\dket\psi,~\hat Z\dket\psi$ with equal prior probabilities of 1/4, where $(\hat X,\hat Y,\hat Z)$ are the three Pauli matrices. This is known as Pauli-twirling. Show that the density matrix received by us is exactly $\hat\rho=\frac12\hat1$. 

Comment: for such a single-qubit mixed state with $\hat\rho=\frac12\hat1$, the outcome for any projective measurement $\hat P$ has a probability of one half, thus the received state contains no information at all about the identity of $\dket\psi$. This is Problem 4 in Chapter 16 exercises of \tong. 



\subsection{Quantum channels}

(1) \textbf{Bonus} Consider an open quantum system $A$ interacting with its environment $B$. In the lectures, we have shown that for an initial state of $\rho_A\otimes\dket{0}\dbra{0}_B$ where $\dket{0}$ is a pure state in $B$, the time evolution for system $A$ is given by a quantum channel $\varepsilon$, which is specified by a set of Kraus operators $\{\hat M_k\}$ satisfying the completeness relation 
\begin{align}
\sum_k\hat M_k^\dagger\hat M_k=\hat 1_A.
\end{align}
Show that this is still true if we consider the initial state to be a product state $\rho_A\otimes\rho_B$, where $\rho_B$ is an arbitrary density matrix in $B$.

Hint: you will need to find the set of Kraus operators satisfying the completeness relation in this case. It is convenient to work in the eigenbasis of $\rho_B$.  
 
(2) \nielsen~Exercise 8.15

%(3) \nielsen~Exercise 8.18

(3) \textbf{Bonus: generalized depolarizing channel} \nielsen~Exercise 8.19

%(5) \nielsen~Exercise 8.28




\subsection{The master equation}

\subsubsection{Phase damping}

Time evolution of a single-qubit open system under the phase damping quantum channel $\varepsilon_z$ is given by:
\begin{align}
\hat\rho=\bpm\rho_{11}&\rho_{12}\\ \rho_{21}&\rho_{22}\epm\rightarrow\varepsilon_z[\hat\rho]=(1-p)\hat\rho+p\hat Z\hat\rho \hat Z=\bpm\rho_{11}&(1-2p)\rho_{12}\\ (1-2p)\rho_{21}&\rho_{22}\epm
\end{align}
where $p$ is the probability of error. For the time evolution by a small time $\Delta t$, the error probability is given by
\begin{align}
2p=\Gamma\Delta t
\end{align}
where $\Gamma$ is the rate of error. 

(1) Derive the differential equation 
\begin{align}
\frac{\text{d}\rho}{\text{d}t}=\lim_{\Delta t\rightarrow0^+}\frac{\varepsilon_z[\rho]-\rho(t=0)}{\Delta t}
\end{align}
and solve it. Obtain the form of $\rho(t)$ for a given initial state $\hat\rho(t=0)=\bpm\rho_{11}&\rho_{12}\\ \rho_{21}&\rho_{22}\epm$. 

(2) Write the differential equation in part (1) into the form of a master equation. Find the Hamiltonian $\hat H$ and the Lindblad operators $\{\hat L_k\}$ in the master equation. 


\subsubsection{Amplitude damping}

(1) Solve the following master equation for a 1-qubit density matrix $\rho(t)$, with a single Lindblad operator $L=\sqrt{\lambda}\sigma_-\equiv\sqrt{\lambda}(\hat X-\imth\hat Y)/2=\sqrt{\lambda}\bpm0&0\\1&0\epm$:
\begin{align}
\frac{\text{d}\rho}{\text{d}t}=\lambda\sigma_-\rho\sigma_+-\frac\lambda2\{\rho,\sigma_+\sigma_-\}
\end{align}
where $\sigma_+=(\hat X+\imth\hat Y)/2=\bpm0&1\\0&0\epm$. Express $\rho(t)$ as a function of the initial density matrix
\begin{align}\label{amplitude damping:initial state}
\rho(t=0)=\bpm\rho_{11}&\rho_{12}\\ \rho_{21}&\rho_{22}\epm
\end{align}

Comment: this master equation describes the amplitude damping of a 1-qubit system. 

(2) Consider the time evolution of a 1-qubit system by an amplitude-damping quantum channel $\varepsilon_p$ with the following set of Kraus operators:
\begin{align}
\hat M_0=\bpm\sqrt{1-p}&0\\0&1\epm,~~~\hat M_1=\bpm0&0\\ \sqrt p&0\epm
\end{align}
For the same initial state (\ref{amplitude damping:initial state}), obtain the final density matrix $\rho(p)=\varepsilon_p[\rho(t=0)]$ after the action of the quantum channel, as a function of the amplitude-damping probability $p$. Compare it with the previous problem, and obtain the relation between $p$ in this problem, and $t$ in the previous problem. 

\subsubsection{Amplitude + phase dampings}

\textbf{Bonus} \nielsen~Exercise 8.30. There is a typo in \nielsen~for this problem: you need to show that $T_2=2T_1$ for amplitude damping only, and $T_2\leq 2T_1$ for both amplitude and phase dampings. 
